\documentclass{article}
\usepackage{graphicx} % Required for inserting images
\usepackage{float}
\usepackage[a4paper, margin=2.5cm]{geometry}

\title{Super Resolution su immagini FFHQ: confronto tra Total Variation e U-Net}
\author{Vincenzo Lombardi}
\date{May 2026}

\begin{document}

\maketitle

\section{Introduzione}

Negli ultimi anni, il \textit{Computational Imaging} ha assunto un ruolo sempre più centrale in numerosi ambiti applicativi, tra cui imaging medico, fotografia computazionale, microscopia e telerilevamento. In questi contesti, l’obiettivo principale consiste nel ricostruire o migliorare un’immagine a partire da osservazioni indirette, incomplete o degradate, sfruttando modelli matematici e algoritmi di ottimizzazione avanzati.

Molti problemi di imaging possono essere formulati come \textit{inverse problems}, nei quali si vuole stimare un’immagine incognita ad alta qualità a partire da una misura osservata affetta da degradazione e rumore. In forma generale, il problema può essere espresso come:

\[
y = Ax + n
\]

dove \(x\) rappresenta l’immagine originale ad alta risoluzione, \(A\) è l’operatore di degradazione (ad esempio downsampling, blur o compressione), \(y\) è l’osservazione disponibile e \(n\) rappresenta il rumore additivo, spesso modellato come rumore gaussiano.

Tra i problemi inversi più rilevanti nel contesto dell’elaborazione di immagini, la \textit{Super Resolution} (SR) riveste particolare importanza. Essa consiste nel ricostruire un’immagine ad alta risoluzione (\textit{High Resolution}, HR) a partire da una versione a bassa risoluzione (\textit{Low Resolution}, LR), cercando di recuperare dettagli fini, texture e strutture che risultano parzialmente o totalmente perse durante il processo di degradazione.

La Super Resolution rappresenta un problema fortemente mal posto (\textit{ill-posed}), poiché molte immagini HR differenti possono generare la stessa osservazione LR. Per questo motivo, la sola inversione dell’operatore di degradazione non è sufficiente e risulta necessario introdurre informazioni a priori sulla soluzione desiderata attraverso tecniche di regolarizzazione.

Tradizionalmente, questo approccio viene affrontato mediante metodi \textit{model-based}, nei quali la soluzione viene ottenuta minimizzando una funzione obiettivo composta da un termine di fedeltà ai dati e un termine di regolarizzazione:

\[
\hat{x} =
\arg\min_x
\left\{
\|Ax - y\|^2 + \lambda R(x)
\right\}
\]

dove \(R(x)\) rappresenta il termine di regolarizzazione e \(\lambda\) ne controlla il peso. Tra le regolarizzazioni più utilizzate vi è la \textit{Total Variation} (TV), particolarmente efficace nel preservare i bordi e ridurre il rumore, pur presentando il limite di produrre talvolta un eccessivo effetto di oversmoothing.

Parallelamente, negli ultimi anni, i metodi basati su \textit{Deep Learning} hanno mostrato risultati molto competitivi nella risoluzione di problemi inversi. In particolare, architetture neurali come la \textit{U-Net} permettono di apprendere direttamente una mappatura tra immagine degradata e immagine ricostruita, sfruttando grandi quantità di dati di training e una struttura encoder-decoder con connessioni skip che favoriscono il recupero dei dettagli spaziali.

In questo progetto viene affrontato il problema della Super Resolution sul dataset FFHQ (\textit{Flickr-Faces-HQ}), costituito da immagini facciali ad alta qualità e risoluzione \(256 \times 256\). L’obiettivo è confrontare due approcci differenti:

\begin{itemize}
    \item un metodo variazionale basato su \textit{Total Variation Regularization};
    \item un metodo \textit{End-to-End} basato su architettura \textit{U-Net}.
\end{itemize}

Le immagini vengono degradate utilizzando due fattori di super-risoluzione (\(\times 2\) e \(\times 4\)) e due livelli di rumore gaussiano (\(\sigma = 0.005\) e \(\sigma = 0.01\)), mantenendo le stesse condizioni sperimentali per entrambi i metodi al fine di garantire un confronto equo.

Le prestazioni vengono valutate sia dal punto di vista quantitativo, mediante le metriche PSNR (\textit{Peak Signal-to-Noise Ratio}) e SSIM (\textit{Structural Similarity Index}), sia dal punto di vista qualitativo attraverso il confronto visivo delle immagini ricostruite.

L’obiettivo finale del lavoro è analizzare vantaggi, limiti e comportamento dei due approcci, evidenziando come, in determinati contesti applicativi, i metodi variazionali classici possano ancora risultare estremamente competitivi rispetto a soluzioni deep learning più complesse.

\newpage
\section{Formulazione del Problema}

Il problema della \textit{Super Resolution} può essere formalizzato come un problema inverso di ricostruzione, nel quale si desidera stimare un’immagine ad alta risoluzione a partire da una sua osservazione degradata a bassa risoluzione.

Dal punto di vista matematico, il modello di acquisizione può essere espresso come:

\[
y = Ax + n
\]

dove:

\begin{itemize}
    \item \(x \in \mathbb{R}^{N}\) rappresenta l’immagine originale ad alta risoluzione (\textit{High Resolution}, HR);
    \item \(y \in \mathbb{R}^{M}\) rappresenta l’immagine osservata a bassa risoluzione (\textit{Low Resolution}, LR), con \(M < N\);
    \item \(A\) è l’operatore di degradazione che modella il processo di downsampling e, eventualmente, ulteriori effetti di blur;
    \item \(n\) rappresenta il rumore additivo, assunto in questo progetto come rumore gaussiano bianco (\textit{Additive White Gaussian Noise}, AWGN).
\end{itemize}

Nel caso specifico di questo lavoro, l’operatore \(A\) consiste principalmente in una riduzione di risoluzione (\textit{downsampling}) con fattori di scala pari a \(\times 2\) e \(\times 4\), seguita dall’aggiunta di rumore gaussiano con deviazione standard pari a:

\[
\sigma \in \{0.005,\; 0.01\}
\]

Successivamente, l’immagine degradata viene riportata alla dimensione originale tramite interpolazione bicubica, in modo da fornire a tutti i metodi confrontati lo stesso input di partenza.

\subsection{Ill-posedness del problema}

La Super Resolution è un problema fortemente mal posto (\textit{ill-posed}), poiché il processo di degradazione comporta una perdita irreversibile di informazione.

In particolare:

\begin{itemize}
    \item più immagini HR differenti possono produrre la stessa osservazione LR;
    \item l’operatore \(A\) non è invertibile in modo stabile;
    \item la presenza di rumore amplifica ulteriormente l’incertezza della ricostruzione.
\end{itemize}

Di conseguenza, una semplice inversione diretta del modello non è sufficiente a garantire una soluzione accurata e stabile.

Per ottenere una stima affidabile dell’immagine originale è necessario introdurre conoscenza a priori sulla struttura desiderata della soluzione, attraverso opportune tecniche di regolarizzazione.

\subsection{Approccio variazionale}

Uno dei modi più classici per affrontare un problema inverso consiste nella formulazione variazionale, nella quale la soluzione viene ottenuta come minimo di una funzione obiettivo composta da due termini:

\[
\hat{x}
=
\arg\min_x
\left\{
\underbrace{\|Ax - y\|^2}_{\text{data fidelity}}
+
\lambda
\underbrace{R(x)}_{\text{regularization}}
\right\}
\]

Il primo termine, detto \textit{data fidelity}, impone coerenza tra la soluzione ricostruita e i dati osservati, penalizzando differenze elevate tra \(Ax\) e \(y\).

Il secondo termine, detto \textit{regularization}, introduce una prior sulla soluzione e consente di stabilizzare il problema, favorendo immagini più realistiche e meno sensibili al rumore.

Il parametro \(\lambda > 0\) controlla il bilanciamento tra fedeltà ai dati e forza della regolarizzazione:

\begin{itemize}
    \item valori piccoli privilegiano l’aderenza ai dati osservati;
    \item valori grandi privilegiano la regolarità della soluzione.
\end{itemize}

Nel presente progetto, il termine di regolarizzazione scelto per il metodo variazionale è la \textit{Total Variation} (TV), descritta nella sezione successiva.

\subsection{Approccio data-driven}

In alternativa ai metodi model-based, i recenti approcci basati su Deep Learning affrontano il problema della Super Resolution come una mappatura diretta tra input degradato e immagine ricostruita.

In questo caso, una rete neurale convoluzionale apprende dai dati una funzione:

\[
f_{\theta}(y) \approx x
\]

dove \(\theta\) rappresenta i parametri della rete ottimizzati durante la fase di training.

Questo approccio non richiede una definizione esplicita del termine di regolarizzazione, poiché la prior viene appresa implicitamente dal dataset di addestramento.

Nel presente lavoro, il metodo data-driven scelto è una rete \textit{U-Net}, descritta in dettaglio nelle sezioni successive.

\newpage
\section{Metodi Utilizzati}

In questo progetto vengono confrontati due approcci differenti per la risoluzione del problema di Super Resolution:

\begin{itemize}
    \item un metodo variazionale classico basato su \textit{Total Variation Regularization};
    \item un metodo \textit{End-to-End} basato su rete neurale convoluzionale di tipo \textit{U-Net}.
\end{itemize}

L’obiettivo del confronto è analizzare il comportamento di un approccio \textit{model-based}, fondato su una formulazione matematica esplicita del problema inverso, rispetto a un approccio \textit{data-driven}, in cui la ricostruzione viene appresa direttamente dai dati.

\subsection{Metodo variazionale: Total Variation Regularization}

La \textit{Total Variation} (TV) rappresenta una delle tecniche di regolarizzazione più utilizzate nei problemi di image restoration e computational imaging.

L’idea principale consiste nel favorire soluzioni che presentino variazioni spaziali limitate, preservando però i bordi netti e le discontinuità strutturali dell’immagine, che spesso corrispondono a informazioni visivamente rilevanti.

Nel caso della Super Resolution, la ricostruzione viene ottenuta risolvendo il seguente problema di ottimizzazione:

\[
\hat{x}
=
\arg\min_x
\left\{
\|Ax - y\|^2
+
\lambda \, TV(x)
\right\}
\]

dove il termine di regolarizzazione TV è definito come:

\[
TV(x)
=
\sum_{i}
\sqrt{
(\nabla_h x_i)^2 +
(\nabla_v x_i)^2
}
\]

con:

\begin{itemize}
    \item \(\nabla_h\): derivata discreta orizzontale;
    \item \(\nabla_v\): derivata discreta verticale;
    \item \(x_i\): intensità del pixel \(i\).
\end{itemize}

Questo termine penalizza forti oscillazioni locali, riducendo il rumore e favorendo immagini più regolari.

\subsubsection{Vantaggi della TV}

I principali vantaggi della regolarizzazione TV sono:

\begin{itemize}
    \item buona robustezza al rumore gaussiano;
    \item preservazione dei bordi e dei contorni;
    \item elevata stabilità numerica;
    \item interpretabilità teorica del modello;
    \item costo computazionale relativamente contenuto.
\end{itemize}

\subsubsection{Limiti della TV}

I principali limiti sono invece:

\begin{itemize}
    \item effetto di \textit{oversmoothing};
    \item perdita di texture fini;
    \item riduzione dei dettagli ad alta frequenza;
    \item possibile effetto \textit{cartoon-like} nelle regioni uniformi.
\end{itemize}

Nel presente lavoro, la TV viene applicata all’immagine degradata e riportata alla risoluzione originale tramite interpolazione bicubica, utilizzando il metodo \texttt{denoise\_tv\_chambolle} disponibile nella libreria \texttt{scikit-image}.

Il parametro di regolarizzazione \(\lambda\) viene scelto empiricamente mediante una procedura di \textit{grid search}, analizzando le metriche PSNR e SSIM.

\subsection{Metodo End-to-End: U-Net}

Come secondo approccio è stata utilizzata una rete neurale convoluzionale di tipo \textit{U-Net}, una delle architetture più diffuse nei problemi di image restoration, segmentation e inverse problems.

La U-Net è una rete di tipo \textit{encoder-decoder} con connessioni \textit{skip connections}, introdotta originariamente per la segmentazione di immagini biomedicali ma oggi largamente impiegata anche in problemi di ricostruzione.

L’idea principale consiste nell’apprendere direttamente una funzione di ricostruzione:

\[
f_{\theta}(y)
\approx x
\]

dove:

\begin{itemize}
    \item \(y\) è l’immagine degradata;
    \item \(x\) è l’immagine ad alta risoluzione desiderata;
    \item \(\theta\) rappresenta i parametri della rete.
\end{itemize}

\subsubsection{Struttura della rete}

La rete U-Net è composta da due parti principali:

\begin{itemize}
    \item \textbf{Encoder}: estrae progressivamente caratteristiche ad alto livello attraverso convoluzioni e downsampling;
    \item \textbf{Decoder}: ricostruisce l’immagine finale attraverso upsampling e convoluzioni.
\end{itemize}

Le \textit{skip connections} collegano direttamente i livelli simmetrici di encoder e decoder, permettendo di preservare informazioni spaziali fini e dettagli locali che altrimenti verrebbero persi durante il downsampling.

Questo aspetto è particolarmente importante nel problema della Super Resolution, dove il recupero di texture e dettagli ad alta frequenza è fondamentale.

\subsubsection{Vantaggi della U-Net}

I principali vantaggi sono:

\begin{itemize}
    \item capacità di apprendere pattern complessi dai dati;
    \item migliore ricostruzione di texture fini;
    \item forte adattabilità al dataset specifico;
    \item ottime prestazioni in molti problemi di image restoration.
\end{itemize}

\subsubsection{Limiti della U-Net}

I principali limiti sono:

\begin{itemize}
    \item necessità di una fase di training costosa;
    \item forte dipendenza dalla qualità del dataset;
    \item maggiore sensibilità agli iperparametri;
    \item minore interpretabilità teorica rispetto ai metodi variazionali.
\end{itemize}

Nel presente progetto, la rete riceve in input l’immagine degradata già riportata alla dimensione originale tramite interpolazione bicubica e produce direttamente la ricostruzione finale ad alta risoluzione.

L’addestramento viene effettuato minimizzando una loss di tipo MSE (\textit{Mean Squared Error}) tra immagine ricostruita e ground truth.

\newpage
\section{Implementazione}

In questa sezione vengono descritte le scelte implementative adottate per la realizzazione del progetto, includendo il dataset utilizzato, la procedura di degradazione delle immagini, la preparazione dei dati e i parametri di training dei due metodi confrontati.

Tutte le implementazioni sono state sviluppate in Python utilizzando Google Colab come ambiente di esecuzione, sfruttando principalmente le librerie \texttt{PyTorch}, \texttt{NumPy}, \texttt{OpenCV}, \texttt{scikit-image} e \texttt{matplotlib}.

\subsection{Dataset}

Per gli esperimenti è stato utilizzato il dataset \textit{FFHQ} (\textit{Flickr-Faces-HQ}), costituito da immagini facciali ad alta qualità, largamente impiegato nei problemi di image restoration e generative modeling.

Per il progetto è stata utilizzata la versione con immagini di dimensione:

\[
256 \times 256
\]

in formato RGB.

Al fine di mantenere il training computazionalmente sostenibile e coerente con la traccia del progetto, è stato selezionato un sottoinsieme di circa:

\[
4000 \text{ immagini}
\]

estratte casualmente dal dataset completo.

\subsection{Suddivisione del dataset}

Il dataset è stato suddiviso in tre sottoinsiemi distinti:

\begin{itemize}
    \item \textbf{Training set}: 3200 immagini
    \item \textbf{Validation set}: 400 immagini
    \item \textbf{Test set}: 400 immagini
\end{itemize}

La suddivisione è stata effettuata utilizzando una procedura randomizzata con \textit{random seed} fissato per garantire la riproducibilità degli esperimenti.

Il training set viene utilizzato per l’addestramento della U-Net, il validation set per il monitoraggio della loss durante il training e il test set per la valutazione finale delle prestazioni.

\subsection{Preprocessing}

Tutte le immagini sono state preprocessate attraverso i seguenti passaggi:

\begin{itemize}
    \item conversione in formato RGB;
    \item ridimensionamento a \(256 \times 256\);
    \item normalizzazione dei valori di intensità nell’intervallo:
\[
[0,1]
\]
\end{itemize}

La normalizzazione risulta particolarmente importante sia per la stabilità numerica dei metodi variazionali sia per il corretto training della rete neurale.

\subsection{Modello di degradazione}

Per simulare il problema di Super Resolution, ciascuna immagine ad alta risoluzione viene degradata secondo il seguente processo:

\begin{enumerate}
    \item downsampling con fattore di scala:
\[
\times 2
\quad \text{oppure} \quad
\times 4
\]

    \item aggiunta di rumore gaussiano bianco (\textit{AWGN}) con deviazione standard:

\[
\sigma
\in
\{0.005,\; 0.01\}
\]

    \item upsampling bicubico fino alla risoluzione originale \(256 \times 256\)
\end{enumerate}

L’upsampling bicubico non rappresenta una ricostruzione finale, ma viene utilizzato come input comune per entrambi i metodi confrontati, in modo da garantire condizioni sperimentali identiche.

Un esempio di degradazione è mostrato in Figura~\ref{fig:degradation}. Le quattro configurazioni testate sono riportate nella tabella seguente:

\begin{figure}[H]
    \centering
    \includegraphics[width=0.95\textwidth]{degradation_x4_noise0.01.png}
    \caption{Esempio di degradazione applicata a un’immagine del test set. A sinistra è mostrata la ground truth, al centro l’input degradato con fattore di super-risoluzione $\times 4$ e rumore gaussiano $\sigma = 0.01$, mentre a destra è riportata la mappa dell’errore assoluto.}
    \label{fig:degradation}
\end{figure}

\begin{center}
\begin{tabular}{cc}
\hline
Fattore SR & Livello di rumore \\
\hline
$\times 2$ & 0.005 \\
$\times 2$ & 0.01 \\
$\times 4$ & 0.005 \\
$\times 4$ & 0.01 \\
\hline
\end{tabular}
\end{center}

\subsection{Implementazione del metodo TV}

Per il metodo variazionale è stata utilizzata la funzione:

\[
\texttt{denoise\_tv\_chambolle}
\]

della libreria \texttt{scikit-image}, che implementa una minimizzazione basata sulla regolarizzazione di Total Variation.

Il parametro principale del metodo è il peso della regolarizzazione:

\[
\lambda
\]

che controlla il bilanciamento tra riduzione del rumore e preservazione dei dettagli.

La scelta di tale parametro è stata effettuata empiricamente mediante una procedura di \textit{grid search} sui valori:

\[
\lambda
\in
\{0.02,\; 0.05,\; 0.08,\; 0.12,\; 0.20\}
\]

Per ogni configurazione di degradazione è stato selezionato il valore che massimizzava le metriche PSNR e SSIM sul validation set.

\subsection{Implementazione della U-Net}

Per il metodo deep learning è stata implementata una U-Net semplificata, composta da:

\begin{itemize}
    \item blocchi convoluzionali encoder;
    \item max pooling per il downsampling;
    \item blocchi decoder con upsampling;
    \item skip connections tra livelli simmetrici;
    \item output finale RGB.
\end{itemize}

\begin{figure}[H]
    \centering
    \includegraphics[width=0.85\textwidth]{comparison_x2_noise0.005.png}
    \caption{Esempio di ricostruzione ottenuta dai diversi metodi nel caso $\times 2$ con rumore gaussiano $\sigma = 0.005$. La figura mostra ground truth, input degradato, baseline bicubica, TV e U-Net, evidenziando il tipo di output prodotto dalla rete rispetto ai metodi di confronto.}
    \label{fig:unet}
\end{figure}

L’input della rete è costituito dall’immagine degradata e bicubicamente interpolata, mentre l’output è la ricostruzione ad alta risoluzione.

\subsection{Parametri di training}

L’addestramento della rete è stato eseguito utilizzando:

\begin{itemize}
    \item \textbf{Loss function}: Mean Squared Error (MSE)
    \item \textbf{Optimizer}: Adam
    \item \textbf{Learning rate}: \(10^{-3}\)
    \item \textbf{Batch size}: 8
    \item \textbf{Epochs}: 3
\end{itemize}

Per ridurre i tempi computazionali compatibilmente con l’ambiente Colab e con la natura individuale del progetto, il training è stato inizialmente eseguito su un sottoinsieme ridotto del training set.

Questa scelta ha permesso di ottenere risultati confrontabili mantenendo tempi di esecuzione sostenibili, pur limitando parzialmente le prestazioni finali della rete rispetto a un training più esteso.

\newpage
\section{Esperimenti e Metriche di Valutazione}

L’obiettivo principale degli esperimenti è confrontare le prestazioni dei due metodi selezionati — Total Variation Regularization e U-Net — nel problema di Super Resolution, analizzando sia la qualità quantitativa delle ricostruzioni sia la qualità percettiva delle immagini ottenute.

Per garantire un confronto corretto, tutti i metodi vengono applicati sugli stessi input degradati, ottenuti secondo il modello descritto nella sezione precedente.

\subsection{Configurazioni sperimentali}

Sono state considerate quattro diverse configurazioni di degradazione, ottenute combinando due fattori di super-risoluzione e due livelli di rumore gaussiano:

\begin{center}
\begin{tabular}{cc}
\hline
Fattore di scala & Deviazione standard del rumore \\
\hline
$\times 2$ & 0.005 \\
$\times 2$ & 0.01 \\
$\times 4$ & 0.005 \\
$\times 4$ & 0.01 \\
\hline
\end{tabular}
\end{center}

Il caso con fattore \(\times 2\) rappresenta una degradazione moderata, mentre il caso \(\times 4\) risulta significativamente più complesso a causa della maggiore perdita di informazione.

Analogamente, il livello di rumore \(\sigma = 0.01\) introduce una perturbazione più severa rispetto a \(\sigma = 0.005\), rendendo il problema maggiormente instabile.

\subsection{Baseline bicubica}

Prima di applicare i metodi di ricostruzione, viene considerata come baseline la semplice interpolazione bicubica.

In questo caso, l’immagine a bassa risoluzione viene riportata alla dimensione originale senza alcuna procedura di regolarizzazione o apprendimento, utilizzando esclusivamente un algoritmo standard di interpolazione.

La baseline bicubica rappresenta un riferimento fondamentale, poiché permette di valutare il reale contributo dei metodi avanzati rispetto a una soluzione semplice e computazionalmente economica.

Il confronto finale viene quindi effettuato tra:

\begin{itemize}
    \item Bicubic interpolation
    \item TV Regularization
    \item U-Net
\end{itemize}

\subsection{Metriche quantitative}

Per la valutazione numerica delle prestazioni sono state utilizzate due metriche standard nei problemi di image restoration:

\begin{itemize}
    \item PSNR (\textit{Peak Signal-to-Noise Ratio})
    \item SSIM (\textit{Structural Similarity Index Measure})
\end{itemize}

\subsubsection{PSNR}

Il PSNR misura la differenza pixel-wise tra immagine ricostruita e ground truth, ed è definito come:

\[
PSNR
=
10 \log_{10}
\left(
\frac{MAX^2}{MSE}
\right)
\]

dove:

\begin{itemize}
    \item \(MAX\) è il valore massimo dell’intensità (nel nostro caso pari a 1);
    \item \(MSE\) è il \textit{Mean Squared Error} tra immagine ricostruita e ground truth.
\end{itemize}

Valori più elevati di PSNR indicano una ricostruzione più accurata.

\subsubsection{SSIM}

Lo SSIM misura la similarità strutturale tra due immagini, considerando luminanza, contrasto e struttura locale.

A differenza del PSNR, lo SSIM è più vicino alla percezione visiva umana e risulta particolarmente utile quando si vogliono valutare qualità percettive e preservazione dei dettagli.

Lo SSIM assume valori nell’intervallo:

\[
0 \leq SSIM \leq 1
\]

dove valori più vicini a 1 indicano una maggiore similarità con la ground truth.

\subsection{Valutazione qualitativa}

Oltre alle metriche numeriche, è stata effettuata anche una valutazione qualitativa mediante confronto visivo diretto tra:

\begin{itemize}
    \item immagine Ground Truth
    \item immagine degradata
    \item ricostruzione Bicubic
    \item ricostruzione TV
    \item ricostruzione U-Net
\end{itemize}

Particolare attenzione è stata posta su:

\begin{itemize}
    \item dettagli del volto;
    \item occhi e sopracciglia;
    \item texture della pelle;
    \item capelli e contorni del viso.
\end{itemize}

In alcune situazioni sono stati utilizzati anche \textit{crop} locali per evidenziare meglio differenze che sull’intera immagine risultavano meno evidenti.

Questa fase è particolarmente importante, poiché metriche molto vicine possono comunque produrre differenze visive significative.

\subsection{Obiettivo dell’analisi}

L’analisi finale non si limita a identificare il metodo con i valori migliori di PSNR e SSIM, ma mira a comprendere:

\begin{itemize}
    \item quando un metodo risulta preferibile rispetto all’altro;
    \item il compromesso tra qualità, stabilità e costo computazionale;
    \item il ruolo della regolarizzazione classica rispetto agli approcci deep learning.
\end{itemize}

Questo confronto permette di valutare criticamente se, nel contesto specifico della Super Resolution su immagini facciali, i metodi variazionali tradizionali possano ancora competere con approcci data-driven più complessi.

\newpage
\section{Risultati Sperimentali}

In questa sezione vengono presentati i risultati ottenuti confrontando i tre approcci considerati:

\begin{itemize}
    \item interpolazione bicubica (baseline);
    \item metodo variazionale con Total Variation (TV);
    \item metodo End-to-End basato su U-Net.
\end{itemize}

L’analisi viene condotta sia dal punto di vista quantitativo, mediante le metriche PSNR e SSIM, sia dal punto di vista qualitativo attraverso il confronto visivo delle immagini ricostruite.

\subsection{Risultati quantitativi}

La Tabella~\ref{tab:results} riporta i valori medi di PSNR e SSIM ottenuti sui campioni del test set per ciascuna configurazione sperimentale.

\begin{figure}[H]
    \centering
    \includegraphics[width=0.95\textwidth]{psnr_summary.png}
    \caption{Confronto dei valori medi di PSNR ottenuti dai diversi metodi sulle quattro configurazioni sperimentali. Il metodo TV ottiene il valore migliore in tutte le configurazioni considerate.}
    \label{fig:psnr_summary}
\end{figure}

\begin{figure}[H]
    \centering
    \includegraphics[width=0.95\textwidth]{ssim_summary.png}
    \caption{Confronto dei valori medi di SSIM ottenuti dai diversi metodi sulle quattro configurazioni sperimentali. Anche in termini di similarità strutturale, la regolarizzazione TV risulta il metodo più stabile e competitivo.}
    \label{fig:ssim_summary}
\end{figure}

\begin{table}[H]
\centering
\caption{Confronto quantitativo tra Bicubic, TV e U-Net}
\label{tab:results}
\begin{tabular}{ccccc}
\hline
Metodo & Scala & Rumore & PSNR & SSIM \\
\hline

Bicubic & $\times 2$ & 0.005 & 32.17 & 0.927 \\
TV & $\times 2$ & 0.005 & \textbf{32.37} & \textbf{0.933} \\
U-Net & $\times 2$ & 0.005 & 32.23 & 0.928 \\

\hline

Bicubic & $\times 2$ & 0.01 & 31.61 & 0.899 \\
TV & $\times 2$ & 0.01 & \textbf{32.17} & \textbf{0.924} \\
U-Net & $\times 2$ & 0.01 & 31.65 & 0.901 \\

\hline

Bicubic & $\times 4$ & 0.005 & 26.26 & 0.800 \\
TV & $\times 4$ & 0.005 & \textbf{26.57} & \textbf{0.802} \\
U-Net & $\times 4$ & 0.005 & 26.50 & 0.802 \\

\hline

Bicubic & $\times 4$ & 0.01 & 26.11 & 0.774 \\
TV & $\times 4$ & 0.01 & \textbf{26.56} & \textbf{0.801} \\
U-Net & $\times 4$ & 0.01 & 26.35 & 0.778 \\

\hline
\end{tabular}
\end{table}

Dall’analisi quantitativa emerge chiaramente che il metodo basato su Total Variation ottiene le migliori prestazioni in tutte le configurazioni testate.

In particolare, il miglioramento risulta più evidente nei casi con rumore più elevato (\(\sigma = 0.01\)), dove la regolarizzazione TV riesce a gestire in modo più efficace la componente rumorosa rispetto sia alla baseline bicubica sia alla rete U-Net.

La U-Net mostra comunque prestazioni competitive e generalmente superiori alla semplice interpolazione bicubica, ma non riesce a superare il metodo variazionale nelle condizioni sperimentali considerate.

\subsection{Analisi del caso \(\times 2\)}

Nel caso di super-risoluzione con fattore \(\times 2\), il problema risulta relativamente meno complesso, poiché la perdita di informazione durante il downsampling è più contenuta.

Con rumore basso (\(\sigma = 0.005\)), tutti i metodi ottengono valori elevati di PSNR e SSIM, con differenze relativamente ridotte.

Quando il rumore aumenta (\(\sigma = 0.01\)), il vantaggio della TV diventa più evidente: il metodo riesce infatti a migliorare sia la qualità numerica sia la stabilità della ricostruzione, mostrando una maggiore robustezza rispetto agli altri approcci.

\subsection{Analisi del caso \(\times 4\)}

Nel caso \(\times 4\), la difficoltà del problema aumenta sensibilmente, poiché il downsampling comporta una perdita molto più severa di dettagli e informazioni ad alta frequenza.

Come previsto, i valori di PSNR e SSIM diminuiscono per tutti i metodi.

Anche in questo scenario, la TV mantiene le migliori prestazioni complessive, mentre la U-Net si avvicina ai risultati del metodo variazionale ma senza riuscire a superarlo.

Questo comportamento conferma come, in presenza di forte degradazione e rumore moderato, la regolarizzazione esplicita possa risultare ancora estremamente efficace.

\subsection{Risultati qualitativi}

Dal punto di vista visivo, le differenze tra i metodi risultano particolarmente evidenti osservando dettagli locali come occhi, sopracciglia, capelli e texture della pelle.

\begin{figure}[H]
    \centering
    \includegraphics[width=\textwidth]{comparison_x4_noise0.01.png}
    \caption{Confronto qualitativo tra ground truth, input degradato, interpolazione bicubica, TV e U-Net nel caso più severo considerato: fattore di super-risoluzione $\times 4$ e rumore gaussiano $\sigma = 0.01$. La TV mostra una maggiore robustezza al rumore, mentre la U-Net conserva risultati competitivi ma più sensibili ai limiti del training.}
    \label{fig:comparison_x4}
\end{figure}

La baseline bicubica produce immagini generalmente più sfocate e meno definite, con una perdita evidente dei dettagli fini.

Il metodo TV riesce a migliorare significativamente la nitidezza dei contorni e a ridurre il rumore, preservando in modo efficace la struttura globale del volto.

Tuttavia, in alcune aree ad alta frequenza può comparire un leggero effetto di \textit{oversmoothing}, che riduce la naturalezza di alcune texture.

La U-Net, invece, tende a preservare meglio alcuni dettagli locali e texture sottili, soprattutto nelle regioni dei capelli e del contorno degli occhi, ma mostra una maggiore sensibilità al rumore residuo e risultati meno stabili tra le diverse configurazioni.

\subsection{Discussione}

Un aspetto particolarmente interessante del progetto è che il metodo deep learning non ha superato il metodo variazionale classico.

Questo risultato non rappresenta un fallimento della U-Net, ma evidenzia piuttosto due aspetti fondamentali:

\begin{enumerate}
    \item la Total Variation rappresenta ancora oggi una regolarizzazione estremamente forte e competitiva nei problemi di image restoration;
    \item le prestazioni delle reti neurali dipendono fortemente dalla quantità di dati, dal tempo di training e dalla scelta degli iperparametri.
\end{enumerate}

Nel presente lavoro, per motivi computazionali, la U-Net è stata addestrata con un numero ridotto di epoche e su un sottoinsieme limitato del dataset.

Un training più esteso avrebbe probabilmente consentito un miglioramento ulteriore delle prestazioni, potenzialmente riducendo il gap rispetto al metodo TV.

Tuttavia, i risultati ottenuti mostrano chiaramente che, in questo specifico scenario, i metodi variazionali classici rimangono altamente competitivi anche rispetto ad approcci deep learning più complessi.

\newpage
\section{Conclusioni}

In questo progetto è stato affrontato il problema della \textit{Super Resolution} su immagini facciali del dataset FFHQ, confrontando due approcci differenti appartenenti a due paradigmi distinti del Computational Imaging:

\begin{itemize}
    \item un metodo variazionale classico basato su \textit{Total Variation Regularization};
    \item un metodo \textit{End-to-End} basato su rete neurale convoluzionale di tipo \textit{U-Net}.
\end{itemize}

Entrambi i metodi sono stati valutati nelle stesse condizioni sperimentali, utilizzando due fattori di super-risoluzione (\(\times 2\) e \(\times 4\)) e due livelli di rumore gaussiano (\(\sigma = 0.005\) e \(\sigma = 0.01\)), con confronto rispetto alla baseline di interpolazione bicubica.

L’analisi quantitativa, effettuata mediante le metriche PSNR e SSIM, ha mostrato che il metodo basato su Total Variation ottiene le migliori prestazioni in tutte le configurazioni testate, risultando particolarmente efficace nei casi con rumore più elevato e nelle situazioni di degradazione più severa.

Dal punto di vista qualitativo, la TV ha dimostrato una buona capacità di preservare i bordi e migliorare la nitidezza complessiva dell’immagine, pur presentando in alcuni casi un leggero effetto di \textit{oversmoothing} nelle texture più fini.

La U-Net ha mostrato risultati competitivi e una buona capacità di recuperare dettagli locali, ma non è riuscita a superare il metodo variazionale nelle condizioni considerate. Questo comportamento è fortemente influenzato dalla limitata fase di training adottata nel progetto, resa necessaria dai vincoli computazionali dell’ambiente Colab e dalla natura individuale del lavoro.

Uno degli aspetti più significativi emersi da questo studio è che i metodi variazionali classici, nonostante la crescente diffusione delle tecniche deep learning, possono ancora rappresentare una soluzione estremamente efficace e robusta nei problemi di image restoration.

In particolare, quando il problema presenta rumore moderato, dataset limitati o risorse computazionali contenute, una buona regolarizzazione esplicita può risultare persino più performante di una rete neurale non completamente ottimizzata.

Questo conferma come, nel Computational Imaging moderno, i metodi \textit{model-based} e quelli \textit{data-driven} non debbano essere considerati necessariamente alternativi, ma spesso complementari.

Come possibili sviluppi futuri del lavoro si possono considerare:

\begin{itemize}
    \item un training più approfondito della U-Net con un numero maggiore di epoche;
    \item l’utilizzo di loss functions percettive (\textit{perceptual loss});
    \item il confronto con metodi ibridi come Plug-and-Play Priors o Diffusion-based restoration;
    \item l’estensione del confronto a dataset non facciali o a problemi di deblurring e denoising più complessi.
\end{itemize}

In conclusione, il progetto ha permesso di evidenziare come, anche in un contesto dominato dal Deep Learning, i metodi variazionali ben progettati rimangano strumenti fondamentali e altamente competitivi nella risoluzione dei problemi inversi.

\newpage
\begin{thebibliography}{9}

\bibitem{rof}
L. Rudin, S. Osher, E. Fatemi,
\textit{Nonlinear total variation based noise removal algorithms},
Physica D, 1992.

\bibitem{unet}
O. Ronneberger, P. Fischer, T. Brox,
\textit{U-Net: Convolutional Networks for Biomedical Image Segmentation},
MICCAI, 2015.

\bibitem{wang2004ssim}
Z. Wang, A. Bovik, H. Sheikh, E. Simoncelli,
\textit{Image quality assessment: From error visibility to structural similarity},
IEEE Transactions on Image Processing, 2004.

\bibitem{sr}
C. Dong, C. C. Loy, K. He, X. Tang,
\textit{Image Super-Resolution Using Deep Convolutional Networks},
IEEE TPAMI, 2016.

\end{thebibliography}

\end{document}
